%%%%%%%%%%%%%%%%%%%%%%%%%%%%%%%%%%%%%%%%%%%%%%%%%%%%%
%%% Task 4 %%%%%%%%%%%%%%%%%%%%%%%%%%%%%%%%%%%%%%%%%%
%%%%%%%%%%%%%%%%%%%%%%%%%%%%%%%%%%%%%%%%%%%%%%%%%%%%%
\task{Emergency power for a hospital}

%%%%%%%%%%%%%%%%%%%%%%%%%%%%%%%%%%%%%%%%%%%%%%
\taskGerman{Notstromversorgung für ein Krankenhaus}

A hospital is upgrading its emergency power system. During normal operation, the critical care units
are supplied from the utility grid. During outages, a diesel generator forms a three-phase line-to-line
\SI{400}{\volt}, \SI{50}{\hertz} AC grid. A dedicated AC/DC converter is used to supply the DC input of a
specific piece of emergency medical equipment. The engineering team is considering the AC/DC converter
topologies shown in Fig. \ref{fig:M3U_vs_M3C} to supply a regulated DC bus.
Assume that the DC load current is continuous and that all semiconductor devices are ideal.
The equipment is approximated by a constant current load.

\begin{germanblock}
	Ein Krankenhaus rüstet seine Notstromversorgung auf. Im Normalbetrieb werden die Intensivstationen aus
	dem öffentlichen Stromnetz versorgt. Bei einem Stromausfall erzeugt ein Dieselgenerator ein Drehstromnetz
	(Leiterspannung \SI{400}{\volt}, \SI{50}{\hertz}). Zur Versorgung medizinischer Notfallgeräte benötigt
	das Krankenhaus einen stabilen Gleichspannungszwischenkreis. Das Entwicklungsteam prüft die in Abb. \ref{fig:M3U_vs_M3C}
	dargestellten AC/DC-Wandlertopologien zur Versorgung eines geregelten Gleichspannungszwischenkreises.
	Es wird angenommen, dass der Gleichstrom lastseitig lückenlos fließt und alle Halbleiterbauelemente ideal
	sind. Die zu versorgende medizinische Ausrüstung wird durch eine konstante Stromlast angenähert.
\end{germanblock}
\begin{figure}[htbp]
	\centering
	\begin{subfigure}[b]{\textwidth}
		\centering
		\begin{circuitikz}
			\def\vd{1cm} % vertical distance inductors
			\def\htraf{0.75cm} % horizontal distance transformer coils
			\draw (0,0) to [short, o-] ++(0.5,0) coordinate (L1astart) to [short] ++(0.5,0) to [L] ++(2,0) coordinate (L1aend)
				(0,-1*\vd) to [short, o-] ++(1,0) coordinate (L1bstart) to [L] ++(2,0) coordinate (L1bend)
				(0,-2*\vd) to [short, o-] ++(1,0) coordinate (L1cstart) to [L] ++(2,0) coordinate (L1cend) -- ++(0,-0.5*\vd) to (\tikztostart -| L1astart)
				to [crossing] ++(0, 1*\vd) to [crossing] ++(0, 1*\vd) to [short, -*] (L1astart)
				(L1aend) -- ++(0,-0.5*\vd) to (\tikztostart -| L1bstart) to [short, -*] (L1bstart)
				(L1bend) -- ++(0,-0.5*\vd) to (\tikztostart -| L1cstart) to [short, -*] (L1cstart);
			\draw let \p1=(L1aend) in (\x1 + \htraf, \y1) coordinate (L2astart) to [L, v^<=$u_{1\mathrm{a}}(t)$, voltage = straight] ++(2,0) to [short, i=$i_{1\mathrm{a}}(t)$] ++(0.5,0) coordinate (L2aend);
			\draw let \p1=(L1bend) in (\x1 + \htraf, \y1) coordinate (L2bstart) to [L, v^<=$u_{1\mathrm{b}}(t)$, voltage = straight] ++(2,0) to [short, i=$i_{1\mathrm{b}}(t)$] ++(0.5,0) coordinate (L2bend);
			\draw let \p1=(L1cend) in (\x1 + \htraf, \y1) coordinate (L2cstart) to [L, v^<=$u_{1\mathrm{c}}(t)$, voltage = straight] ++(2,0) to [short, i=$i_{1\mathrm{c}}(t)$] ++(0.5,0)  coordinate (L2cend);
			\draw (L2astart) to [short, -*] (L2bstart) to [short, -*] (L2cstart) -- ++(0, -1*\vd) -- ++(5,0) coordinate (Rend);
			\draw[double, double distance=3pt, thick] let \p1=(L1aend), \p2=(L2cstart) in (\x1/2+\x2/2, \y1) -- (\x1/2+\x2/2, \y2);
			\draw (L2aend) to [diode] ++(1.25,0) coordinate (D1end);
			\draw (L2bend) to [diode] ++(1.25,0) coordinate (D2end);
			\draw (L2cend) to [diode] ++(1.25,0) coordinate (D3end) to [short, -*] (D2end) to [short, -*] (D1end);
			\draw (D1end) to [short] ++(0.5,0) coordinate (u2) to [short, i=$i_2(t)$] ++(2.0,0) to [I] (Rend -| \tikztostart) to (Rend);
			\draw (u2) to [open, v^>=$u_2(t)$, voltage shift=2pt, voltage = straight] (Rend -| \tikztostart);
		\end{circuitikz}
		\caption{M3U: Three-phase three-pulse uncontrolled mid-point rectifier.}
		\label{fig:M3U}
	\end{subfigure}
	\hfill
	\begin{subfigure}[b]{\textwidth}
		\centering
		\begin{circuitikz}
			\def\vd{1cm} % vertical distance inductors
			\def\htraf{0.75cm} % horizontal distance transformer coils
			\draw (0,0) to [short, o-] ++(0.5,0) coordinate (L1astart) to [short] ++(0.5,0) to [L] ++(2,0) coordinate (L1aend)
				(0,-1*\vd) to [short, o-] ++(1,0) coordinate (L1bstart) to [L] ++(2,0) coordinate (L1bend)
				(0,-2*\vd) to [short, o-] ++(1,0) coordinate (L1cstart) to [L] ++(2,0) coordinate (L1cend) -- ++(0,-0.5*\vd) to (\tikztostart -| L1astart)
				to [crossing] ++(0, 1*\vd) to [crossing] ++(0, 1*\vd) to [short, -*] (L1astart)
				(L1aend) -- ++(0,-0.5*\vd) to (\tikztostart -| L1bstart) to [short, -*] (L1bstart)
				(L1bend) -- ++(0,-0.5*\vd) to (\tikztostart -| L1cstart) to [short, -*] (L1cstart);
			\draw let \p1=(L1aend) in (\x1 + \htraf, \y1) coordinate (L2astart) to [L, v^<=$u_{1\mathrm{a}}(t)$, voltage = straight] ++(2,0) to [short, i=$i_{1\mathrm{a}}(t)$] ++(0.5,0) coordinate (L2aend);
			\draw let \p1=(L1bend) in (\x1 + \htraf, \y1) coordinate (L2bstart) to [L, v^<=$u_{1\mathrm{b}}(t)$, voltage = straight] ++(2,0) to [short, i=$i_{1\mathrm{b}}(t)$] ++(0.5,0) coordinate (L2bend);
			\draw let \p1=(L1cend) in (\x1 + \htraf, \y1) coordinate (L2cstart) to [L, v^<=$u_{1\mathrm{c}}(t)$, voltage = straight] ++(2,0) to [short, i=$i_{1\mathrm{c}}(t)$] ++(0.5,0)  coordinate (L2cend);
			\draw (L2astart) to [short, -*] (L2bstart) to [short, -*] (L2cstart) -- ++(0, -1*\vd) -- ++(5,0) coordinate (Rend);
			\draw[double, double distance=3pt, thick] let \p1=(L1aend), \p2=(L2cstart) in (\x1/2+\x2/2, \y1) -- (\x1/2+\x2/2, \y2);
			\draw (L2aend) to [thyristor] ++(1.25,0) coordinate (D1end);
			\draw (L2bend) to [thyristor] ++(1.25,0) coordinate (D2end);
			\draw (L2cend) to [thyristor] ++(1.25,0) coordinate (D3end) to [short, -*] (D2end) to [short, -*] (D1end);
			\draw (D1end) to [short] ++(0.5,0) coordinate (u2) to [short, i=$i_2(t)$] ++(2.0,0) to [I] (Rend -| \tikztostart) to (Rend);
			\draw (u2) to [open, v^>=$u_2(t)$, voltage shift=2pt, voltage = straight] (Rend -| \tikztostart);
		\end{circuitikz}%
		\caption{M3C: Three-phase three-pulse fully controlled mid-point rectifier.}
		\label{fig:M3C}
	\end{subfigure}
	\caption{Considered three-phase rectifiers. The two input transformers are assumed to be ideal and have both the equal number of turns in the primary and secondary windings.}
	\label{fig:M3U_vs_M3C}
\end{figure}
\subtask{During emergency operation, the generator terminal voltage may vary by $\pm \SI{15}{\percent}$
	due to operating conditions. Which converter topology (M3U or M3C) is more suitable for maintaining a desired
	DC output voltage?}{2}

%
\subtaskGerman{Im Notbetrieb kann die Klemmenspannung des Generators aufgrund der Betriebsbedingungen um $\pm \SI{15}{\percent}$
	schwanken. Welche Umrichtertopologie (M3U oder M3C) eignet sich besser, um eine gewünschte Gleichausgangsspannung aufrechtzuerhalten?}
\begin{solutionblock}
	The preferred converter is the three-phase controlled mid-point rectifier: the M3C uses thyristors whose firing angle ($\alpha$) can be adjusted, allowing
	control of the average DC output voltage. If the generator voltage changes, the output voltage can be adjusted
	by changing the firing angle (within the converter's operating limits). In contrast, the M3U cannot actively control the output voltage, which might be problematic when supplying critical medical equipment.
\end{solutionblock}
%
\subtask{Assuming the team selected the M3C converter and initially operates it with a firing angle of $\alpha=\SI{25}{\degree}$. Calculate the average DC output voltage.}{2}
\subtaskGerman{Angenommen, das Team hat sich für den M3C-Stromrichter entschieden und betreibt diesen zunächst
	mit einem Zündwinkel von $\alpha=\SI{25}{\degree}$. Berechnen Sie die mittlere Gleichausgangsspannung.}
\begin{solutionblock}
	For a three-phase fully controlled mid-point rectifier,
	$$\hat{u}_\mathrm{1,normal} = \sqrt{2}\frac{U_{1\mathrm{LL}}}{\sqrt{3}}
		= \sqrt{2}\frac{\SI{400}{\volt}}{\sqrt{3}} = \SI{326.6}{\volt},$$
	the average DC output voltage is given by:
	$$U_\mathrm{2,normal}=\frac{3\sqrt{3}}{2\pi}\hat{u}_\mathrm{1,normal}\cos{\alpha}
		= \SI{244.79}{\volt}.$$
\end{solutionblock}
%
\subtask{During an emergency, the generator terminal voltage suddenly drops by $\SI{15}{\percent}$.
	\begin{enumerate}
		\renewcommand{\labelenumi}{\alph{enumi})}
		\item Calculate the new average DC output voltage, while the firing angle remains unchanged.
		\item Can the original DC output voltage be restored only by changing the firing angle?
		\item Suggest an alternative converter/topology to restore the original output voltage in
		      such emergency cases, and justify your answer with calculations.
	\end{enumerate}}{3}
\subtaskGerman{Im Notfall sinkt die Klemmenspannung des Generators schlagartig um $\SI{15}{\percent}$.
	\begin{enumerate}
		\renewcommand{\labelenumi}{\alph{enumi})}
		\item Berechnen Sie die neue mittlere Gleichausgangsspannung, während der Zündwinkel unverändert bleibt.
		\item Kann die ursprüngliche Gleichausgangsspannung nur durch Änderung des Zündwinkels wiederhergestellt werden?
		\item Schlagen Sie einen alternativen Wandler bzw. eine alternative Topologie vor, um die ursprüngliche
		      Ausgangsspannung in solchen Notfällen wiederherzustellen, und begründen Sie Ihre Antwort mit
		      Berechnungen.
	\end{enumerate}}
\begin{solutionblock}
	\begin{enumerate}
		\renewcommand{\labelenumi}{\alph{enumi})}
		\item $$U_{2,\mathrm{reduced}}=(1-0.15) \cdot U_\mathrm{2,normal} = (1-0.15) \cdot \SI{245}{\volt} = \SI{208.3}{\volt}.$$
		\item The new peak phase input voltage is:
		      $$\hat{u}_{1,\mathrm{reduced}} = (1-0.15) \cdot \hat{u}_\mathrm{1,normal} = (1-0.15)\cdot \SI{326.6}{\volt} = \SI{277.6}{\volt},$$
		      The maximum DC voltage obtainable occurs at $\alpha=\SI{0}{\degree}$. Therefore, the original DC output voltage cannot be restored only by changing the firing angle:
		      $$U_{2,\mathrm{max}} = \frac{3\sqrt{3}}{2\pi}\hat{u}_{1,\mathrm{reduced}}
			      = \SI{229.58}{\volt} < \SI{244.79}{\volt}.$$
		\item B6C (six-pulse fully controlled bridge rectifier) can be used to restore the original output voltage.
		      $$U_2 = \hat{u}_{1,\mathrm{LL},\mathrm{reduced}}\frac{6}{\pi}\sin{\frac{\pi}{6}}\cos{\alpha},$$
		      $$\Leftrightarrow \quad \SI{244.79}{\volt} = \sqrt{2}\cdot0.85\cdot\SI{400}{\volt}\cdot\frac{6}{\pi}\sin{\frac{\pi}{6}}\cos{\alpha},$$
		      $$\Leftrightarrow \quad \alpha = \arccos\left(\frac{\SI{244.79}{\volt}}{\SI{459.16}{\volt}}\right) = \SI{57.8}{\degree}.$$
	\end{enumerate}
\end{solutionblock}

